\section{Увод}
\label{sec:uvod}

По делата за престъпления от общ характер наказателният процес в Република България
преминава през две основни фази~-- досъдебна и съдебна.\footnote{По делата за
престъпления, които се преследват по тъжба на пострадалия, досъдебно производство няма:
делото се образува направо пред съда, а обвинението се поддържа от пострадалия като
частен тъжител (чл.~80, 81 и чл.~247, ал.~1, т.~2 НПК).} В досъдебната фаза се изяснява дали е извършено престъпление и кой го е
извършил, събират се доказателствата и се формулира обвинението, което съдът ще
разглежда. Според Конституцията и Наказателно-процесуалния кодекс тази фаза се ръководи
от прокурора (чл.~127, т.~1--3 КРБ; чл.~46 НПК).\footnote{Конституция на Република
България, обн. ДВ, бр.~56 от 13.07.1991~г. Правомощията на прокуратурата да ръководи
разследването и сама да разследва (чл.~127, т.~1 и 2) са добавени с изменението, обн. ДВ,
бр.~27 от 2006~г.}

Темата е актуална по три причини. Първо, правомощията на прокурора бяха изменяни
неколкократно: през 2023~г. беше въведен съдебен контрол върху отказа да се образува
досъдебно производство за тежки престъпления, а през 2025~г. бяха променени сроковете на
разследването и беше уредено изземването на криптоактиви. Второ, при киберпрестъпленията
доказателствата са електронни, нетрайни и често се намират в чужбина, което поставя пред
прокурора задачи, за които класическите процесуални институти не са създадени. Трето, от
18~август 2026~г. се прилага Регламент (ЕС) 2023/1543 за електронните доказателства.

\textit{Обект} на изследването са отношенията между прокурора, разследващите органи и
съда в досъдебното производство, а \textit{предмет}~-- правомощията на прокурора в тази
фаза по действащото право и упражняването им при киберпрестъпленията. \textit{Целта} е да
се изясни ролята на прокурора в досъдебното производство и да се покаже как тя се
проявява при киберпрестъпленията. \textit{Задачите} са: да се изяснят основите на
прокурорските правомощия; да се проследят правомощията от образуването до приключването
на производството; да се очертае границата между правомощията, които прокурорът
упражнява сам, и тези, които изискват разрешение на съда; да се анализират особеностите
при електронните доказателства и при международното сътрудничество; да се формулират
изводи.

Използвани са нормативният метод, систематичното тълкуване и анализът на практиката на
Съда на ЕС и на Европейския съд по правата на човека. Работата не разглежда участието на
прокурора в съдебната фаза.
